Here we follow the discussion in
\cite{eberhardtOffshellPartitionFunctions2022} to give an
illustration of the modern perspective of the gravitational phase space of AdS3.

\subsection{Naive Phase Space}

We start with the double Chern-Simons Theory to represent the pure
AdS3 gravity as discussed in \cref{sec:AdSCS}. Modernly, we will
slightly modify the theory by using $ PSL(2,\mathbb{R}) $ gauge
group. Locally $ PSL(2,\mathbb{R}) $ is the same as $
SL(2,\mathbb{R}) $, they have the same Lie Algebra
but globally they are different. In our context the difference will
lead to a shift in $ k_{SL} = l/4G $ to $ k_{PSL} = l/16G $, which is
minor in the phase space analysis.

We know that the EoM of Chern-Simons Theory phase space is given by
the space of flat $ PSL(2,\mathbb{R}) $ connection. If we consider
the topology of the spacetime to be $ \mathbb{R} \times \Sigma $,
then the phase space is given by the moduli space of flat $
PSL(2,\mathbb{R}) $ connection on $ \Sigma $:
\begin{align}
  \mathcal{M}_{PSL(2,\mathbb{R})} &= \frac{\text{Flat } PSL(2,\mathbb{R}) \text{
  connections on } \Sigma}{\text{Gauge Transformations}} \\
  &= Hom(\pi_1(\Sigma), PSL(2,\mathbb{R}))/PSL(2,\mathbb{R})
\end{align}
Naively we may say that the gravitational phase space of AdS3 is given by:
\begin{align}
  \mathcal{M}_{grav} \cong \mathcal{M}_{PSL(2,\mathbb{R})} \times
  \mathcal{M}_{PSL(2,\mathbb{R})}
\end{align}
However, one should remeber that the double Chern-Simons Theory is
not exactly equivalent to the pure AdS3 gravity, as we have discussed
in \cref{sec:comparephase}, we have two subtleties left:
\begin{enumerate}
  \item The metric should be non-degenrate or the Vielbein should be
    invertible, which is not guaranteed by the
    double Chern-Simons Theory.
  \item The large diffeomorphisms are not included in the gauge
    transformations of the Chern-Simons Theory, thus we need to mod
    out the large
    diffeomorphisms in the phase space of the double Chern-Simons
    Theory to get the correct gravitational phase space.
\end{enumerate}
But these problems can be solvable in AdS case.

\subsection{Phase Space as Double Teichmüller Space}

\subsubsection{Non-Degenrate Metric and Teichmüller Component}

Just as we discussed in \cref{sec:Topologicalsl2r}, similar to the
case of $ SL(2,\mathbb{R}) $, the moduli space of flat
$PSL(2,\mathbb{R}) $ connection on $ \Sigma $ is given by several
disconnected components, which are labeled by the Euler class of the
flat plane bundle.

Given by paper \cite{krasnovMinimalSurfacesParticles2007,
scarinciUniversalPhaseSpace2013} we have an important geometric theorem:
\thm{
  In the case of AdS3, the non-degenrate metric condition is
  equivalent to taking the $ e = 2g -2 $ component of the moduli
  space of flat $PSL(2,\mathbb{R}) $ connection on $ \Sigma $.
}
Thus, we only have to focus on the $ e = 2g -2 $ component, which is
also called the
Teichmüller component, since it is isomorphic to the Teichmüller
space of $ \Sigma $, then we might see that the gravitational phase
space of AdS3 is given by the Teichmüller space of $ \Sigma $:
\begin{align}
  \mathcal{M}_{grav} \subset \mathcal{M}_{PSL(2,\mathbb{R})}^{(e = 2g
  -2)} \times \mathcal{M}_{PSL(2,\mathbb{R})}^{(e = 2g -2)} \cong
  \mathcal{T}(\Sigma) \times \mathcal{T}(\Sigma)
\end{align}

\subsubsection{Large Diffeomorphisms and Gauging Mapping Class Group}

Then we have to gauge the large diffeomorphisms. To do this, we note
that the difference between large diffeomorphism and small
diffeomorphism is given by the mapping class group of $ M $:
\begin{align}
  \text{MCG}(M) = \text{Diff}(M)/\text{Diff}_0(M)
\end{align}
We notice that the mapping class group of $ M = \mathbb{R} \times
\Sigma $ is given by the mapping class group of $ \Sigma $:
\begin{align}
  \text{MCG}(M) \cong \text{MCG}(\Sigma)
\end{align}
Thus we can conclude the gravitational phase space of AdS3 is given by:
\thm{
  AdS3 Phase Space

  The gravitational phase space of pure AdS3 gravity is given by:
  \begin{align}
    \mathcal{M}_{grav} = \frac{\mathcal{T}(\Sigma) \times
    \mathcal{T}(\Sigma)}{\text{MCG}(\Sigma)}
  \end{align}
  We need to mode out the \textbf{Diagonal} mapping class group that
  act on both copy of the Teichmüller space, since the large
  diffeomorphisms act on both copy of the Teichmüller space in the same way.
}
In modern context of quantizing the theory, we will mode out the
mapping class group first but quantize the Teichmuller space first
and then mode out the mapping class group when computing the
Partition function in \cite{collierSolving3dGravity2023}.
